%% kapa: exam=1 round=19 year=2008 period=1 form=A subject=civil_law subject_ko=민법 seq=17 number=17
%% source: Q-Net 감정평가사 기출문제 (cst00309) artlSeq=1001962 file=extracted/1008/1001962/2008년 제19회 감정평가사 1차시험 기출문제/1교시.hwp
%% conversion: kordoc 4.24.51 -> tools/gen1.py (auto); review=proofread-sonnet; flags=none
\begin{question}{17}
조건과 기한에 관한 설명으로 틀린 것은? (다툼이 있는 경우에는 판례에 의함)
\begin{choices}
  \item 부관이 붙은 법률행위에서 부관에 표시된 사실이 발생하지 않으면 채무를 이행하지 않아도 된다고 보는 것이 상당한 경우에 그 부관은 조건이 아니라 불확정기한에 해당한다.
  \item 건축허가를 받지 못한 때에는 토지매매계약을 무효로 하기로 한 약정은 해제조건부 법률행위에 해당한다.
  \item 정지조건부 법률행위에서 조건이 성취되었다는 사실은 이에 의하여 권리를 취득하고자 하는 자에게 증명책임이 있다.
  \item 불확정한 사실이 발생한 때를 이행기한으로 정한 경우에는 그 사실이 발생한 때는 물론 그 사실의 발생이 불가능하게 된 때에도 이행기한은 도래한 것으로 보아야 한다.
  \item 조건의 성취가 미정한 권리의무는 일반규정에 의하여 처분, 상속, 보존 또는 담보로 할 수 있다.
\end{choices}
\end{question}
\answer{1}
